\color{unimain}\section[Использование по назначению]{Использование по назначению}\label{sec:using}
\color{black}

\color{unimain}\subsection{Эксплуатационные ограничения}\color{black}
Климатические условия эксплуатации должны соответствовать требованиям документа ЮТКБ.656122.609 РЭ <<Устройство защиты и автоматики серии <<ЮНИТ-М300>>. Руководство по эксплуатации общее>>.

\color{unimain}\subsection{Подготовка устройства к использованию}\color{black}
Меры безопасности при подготовке устройства, порядок проведения внешнего осмотра устройства перед использованием, требования к монтажу устройства, подключение к персональному компьютеру и устройству <<ЮНИТ-ИЧМ>> приводится в документации ЮТКБ.656122.609 РЭ <<Устройство защиты и автоматики серии <<ЮНИТ-М300>>. Руководство по эксплуатации общее>>. Возможность работы устройства в условиях, отличных от указанных, должна согласовываться с предприятием-изготовителем.

%%%%%%%%%%%%%%%%%%%%%%%%%%%%%%%%%%%%%%%%% Работа с устройством %%%%%%%%%%%%%%%%%%%%%%%%%%%%%%%%%%%%%%%%%%

\color{unimain}\subsection{Работа с устройством}\color{black}

\begin{enumerate}

\item Взаимодействие с устройством осуществляется при помощи подключаемого устройства <<ЮНИТ-ИЧМ>> или через персональный компьютер с установленным программным обеспечением <<ЮНИТ-СЕРВИС>>. Руководство по устройству <<ЮНИТ-ИЧМ>> представлено в документе ЮТКБ.656132.701 РЭ <<Блок интерфейсный <<ЮНИТ-ИЧМ>>. Руководство по эксплуатации>>. Руководство по использованию программного обеспечения представлено в документе RU.37182817.00001.02.34.01 <<Программное обеспечение мониторинга и параметрирования <<ЮНИТ-СЕРВИС>>. Руководство пользователя>>.

\item
Описание работы с устройством (включение, способы управления устройством, обновление встроенного программного обеспечения, неисправности и методы их устранения) приводится в документации ЮТКБ.656122.609 РЭ <<Устройство защиты и автоматики серии <<ЮНИТ-М300>>. Руководство по эксплуатации общее>>.

\end{enumerate}

%%%%%%%%%%%%%%%%%%%%%%%%%%%%%%%%%%%%%%%%% Общие рекомендации по выбору параметров функционирования %%%%%%%%%%%%%%%%%%%%%%%%%%%%%%%%%%%%%%%%%%

\color{unimain}\subsection{Общие рекомендации по выбору параметров функционирования}\color{black}

\begin{enumerate}

\item Настройка  устройства  сводится  к  заданию  значений  параметров  для  каждого  конкретного случая применения. Значения параметров конфигурации аналоговых входов, такие как коэффициенты трансформации, фазы подключаемых к устройству цепей переменного тока и виды схем подключения цепей ТН и др., необходимо задать в точном соответствии с вариантом использования устройства.

\item При использовании в устройстве интерфейсов токовой петли 4(0)…20 мА также необходимо задать параметры конфигурации для расчета текущего положения РПН.

\item Режим выдачи команд («Импульсный» или «Непрерывный») устанавливается в соответствие с используемым для совместной работы типом привода. По умолчанию (как наиболее вероятный случай) установлен «Импульсный» режим выдачи команд. 

\item Уставка блокировки регулирования по току должна быть выбрана в соответствие с номинальным током устройства РПН. Уставки номинального тока ввода 1 и 2 секции нужно задать равными номинальному току соответствующей обмотки трансформатора. 

\item
Желаемый  уровень  (или  набор  уровней  для  различных  ситуаций)  поддерживаемого  значения напряжения  в  автоматическом  режиме  должен  определяться  местным  диспетчерским  центром.  При частых  колебаниях  напряжения  в  течение  суточного  графика  нагрузки  устройство,  работающее  в автоматическом  режиме  управления,  приводит  в  действие  устройство  РПН,  что  способствуют  его ускоренному  износу. Поэтому выбор между ручным и автоматическим режимом управления в условиях частых колебаний основывается на балансе необходимости точного поддержания уровня напряжения у потребителя и снижения затрат на ремонт и техническое обслуживание оборудования. В пользу выбора ручного  режима  управления  влияют  такие  факторы,  как  наличие  дежурного  персонала  и  возможность дистанционного управления.

\item
Ширина  зоны  нечувствительности  (в  режиме  автоматического  регулирования)  должна соответствовать ступени регулирования с коэффициентом отстройки, равным 1,4…2,0. Например, если в 
паспортных  данных    трансформатора  указан  диапазон  регулирования  9  +  1,78\%,  параметр, определяющий ширину зоны нечувствительности можно принять равным $\frac{1,7 · 1,78\%}{100\%} = 0,03$.

\item
Время  выдачи  первой  команды  регулирования  («Прибавить»  или  «Убавить»)  в  автоматическом режиме (T1) должно быть отстроено от кратковременных просадок (скачков) напряжения при включении (отключении) нагрузки. Назначение данной выдержки времени – снижение количества коммутаций из-за непродолжительных колебаний напряжения. Обычно принимается равным 150…180 с.

\item
Время выдачи повторной команды (T2) автоматического регулирования можно принять равным 10 с.

\item
То же значение (10 с) можно рекомендовать для задержки (T3) автоматической команды «Убавить» при значительном увеличении напряжения на регулируемой секции.

\item
Время  реакции  привода  (t1)  на  команду  управления  (только  для  импульсного  режима  выдачи управляющих команд) должна быть не более 2 с.

\item
Время  максимальной продолжительности переключения (t2) для импульсного режима управления должно быть заведомо больше длительности цикла переключения привода. 

\item
Задержка  снятия  сигналов  управления  (t3)  после  формирования  приводом  сигнала  «Идет переключение» для импульсного режима может быть принята 0,5 с. 

\end{enumerate}


%\textcolor{red}{Здесь указать какие текущие значения аналоговых величин и где посмотреть. А также наверно либо здесь картинка с Деревом меню либо ссылка на приложение Дерево меню }


